% !TEX root = ../monograph.tex
\section{Scaling dimension and the physical quotient}
\label{ch18:graded-quotient-section}

An estimate on a polynomial coefficient is useful for local product replacement only if
it retains the degree of the field that coefficient represents.
We derive the required quotient estimate from the triangular marked
transfer. This step uses the corrected physical inclusions and the
normalized-null statement of Section~\ref{ch13:quotient}; those statements
were established using the exact microscopic equation space.

Fix an output scaling dimension $t$. For $v\in V_t$, let
$x_l(v)=(x_{l,\alpha}(v))_\alpha$ be the vector of polynomial coefficients
after transport to level $l$, in the fixed basis of polynomials of
scaling dimension at most $t$. Write $|x|_1=\sum_\alpha|x_\alpha|$. The leading polynomial map at one step has the form
$M_l\mathsf d_l$, where
\[
 (\mathsf d_lx)_\alpha=L_l^{t-|\alpha|}x_\alpha.
\]
Here $|\alpha|$ is the scaling dimension of the basis polynomial, not
its factor count. The leading map fixes the constant and is triangular:
$(M_l)_{\alpha\beta}=0$ off the diagonal whenever
$|\alpha|\le|\beta|$. Both $M_l$ and its inverse have norm at most the
prescribed stage constant $K_l$. The full map differs from it by the
marked remainder bound of Section~\ref{ch11:one-graph}.
Enlarge the prescribed column bounds $C_l$ and block floors so that
\begin{equation}
 64K_l^2\le C_l,C_{l+1},\qquad C_{l+1}\le L_l,
 \qquad 2K_lL_l^{-3/8}(1+C_{l+1})\le1.
 \label{ch18:graded-room}
\end{equation}
These are finite-stage inequalities, with constants fixed before the
block ratio, and the schedule gives
$\sum_{l<m}\log C_{l+1}=o(\tau_m)$.

\begin{proposition}[Graded quotient estimate]
\label{ch18:graded-estimate}
For $0\le i\le t$ there are constants $C_{t,i}(m)$ with
$\log C_{t,i}(m)=o(\tau_m)$ such that the following holds.
Suppose $v_{N,p}\in V_t$ converges on an ordered cutoff and volume
refinement to $v$, and $|x_m(v_{N,p})|_1\le\epsilon$ on that tail.
Then
\begin{equation}
 \operatorname{dist}_{F_t}([v],\iota_{i-1,t}F_{i-1})
       \le C_{t,i}(m)a_m^{-(t-i)}\epsilon,
 \label{ch18:graded-distance}
\end{equation}
where $F_{-1}=0$. Constants are independent of the refinement and the
convergent label family. The norm on $F_t$ is fixed once.
\end{proposition}
\begin{proof}
For $i>0$ let $P$ project onto scaling dimensions below $i$.
Triangularity makes $\ker P$ invariant under $M_l$ and its inverse.
The degree scaling expands low coordinates by at least
$L_l^{t-i+1}$, and high coordinates by at most $L_l^{t-i}$.
The one-step marked bound and \eqref{ch18:graded-room} imply
\begin{align}
 |x_{l+1}|_1\le C_{l+1}|Px_{l+1}|_1
 &\ \Longrightarrow\quad
 |x_l|_1\le C_l|Px_l|_1,\quad
 L_l^{t-i+1}|Px_{l+1}|_1\le C_{l+1}|Px_l|_1,
 \label{ch18:low-cone}\\
 Px_{l+1}=0
 &\ \Longrightarrow\quad
 |x_l|_1\le C_{l+1}L_l^{t-i}|x_{l+1}|_1.
 \label{ch18:high-step}
\end{align}
To check the first implication, put $Y=|P\mathsf d_lx_{l+1}|_1$.
The premise and $C_{l+1}\le L_l$ give
$|\mathsf d_lx_{l+1}|_1\le2Y$ and $|x_{l+1}|_1\le Y$. Applying the inverse leading map
to the low part and absorbing the marked error gives
$Y\le2K_l|Px_l|_1$. The full output is at most $3K_lY$;
\eqref{ch18:graded-room} gives both assertions. When the low part
is zero, the high-degree scaling and $K_l+1\le C_{l+1}$ give
\eqref{ch18:high-step}.

At the microscopic level take the low plane $PV_t$ and transport it.
Equation~\eqref{ch18:low-cone} shows that its projection onto the low
coordinates at every later level is injective, and hence onto that
fixed-dimensional coordinate space. At level $m$ subtract its vector
with the same low coordinates as $x_m(v_{N,p})$. The remaining
coefficient has zero low part and norm at most $(1+C_m)\epsilon$.
Repeat the subtraction after each step towards the fixed reference
level, applying \eqref{ch18:high-step} between subtractions. The sum of
all subtracted vectors still lies in the transported microscopic low
plane. Calling its reference label $y_{N,p}$, we obtain
\[
 |v_{N,p}-y_{N,p}|_1\le
 G_m(1+C_m)(a_{r(t)}/a_m)^{t-i}\epsilon,
 \qquad G_m=\prod_{l<m}(1+C_l)C_{l+1}.
\]
Finite initial transfer and norm-equivalence constants are absorbed in
$G_m$. Its logarithm is $o(\tau_m)$ by the prescribed column bounds.

It remains to identify this microscopic low plane with physical lower
degrees; it is not the naive inclusion of unnormalized polynomials.
The labels $y_{N,p}$ are bounded by the displayed inequality.
The expansion in \eqref{ch18:low-cone}, iterated from the microscopic
level, bounds their microscopic low coefficients by
$a_N^{t-i+1}e^{o(\tau_N)}$. Therefore their degree-$t$ physical
representatives can be written with the normalization $a_N^{-(i-1)}$
and coefficient/derivative growth $e^{o(\tau_N)}$, on the same fixed
stencil. Their global value bounds also have this growth, since the
threshold inverses have subpower growth.

The local-trial comparison of Chapter~\ref{ch11:insertions} matches
these representatives to degree-$(i-1)$ reference labels with vanishing
separated residuals. The matched coefficients need not be bounded.
Apply the normalized-null representative argument of
Section~\ref{ch13:quotient}: division by one plus the norm of their
corrected projection makes the equation remainder uniformly controlled,
and its separated responses vanish. Since $y_{N,p}$ has a bounded
reference limit, that argument gives a physical degree-$(i-1)$ field
with exactly its limiting responses. The all-spectator definition and
corrected inclusion thus put every such limit $[y]$ in
$\iota_{i-1,t}F_{i-1}$. Passing to a convergent subsequence of the bounded
$y_{N,p}$ proves \eqref{ch18:graded-distance}. Extending a tail-defined
bounded family by zero outside its tail changes none of these limits.

For $i=0$ no low plane is removed. The reverse coefficient estimate
directly bounds $|v|_1$ by
$e^{o(\tau_m)}(a_{r(t)}/a_m)^t\epsilon$. Passing to the quotient
proves the same conclusion. All fixed reference spacings are absorbed
in $C_{t,i}(m)$.
\end{proof}

The power $t-i$ in \eqref{ch18:graded-distance} is what recovers a
coefficient of degree $i$ from an auxiliary extraction at a much higher
degree $t$. An ungraded bound on $[v]$ would lose precisely that power.
The proof also explains why identifying the lower field requires all
spectator responses, even though its eventual coefficient is measured
by a fixed finite family of probes.
